\chapter{相关工作}

\label{cha:ralated work}


\section{WebAssembly 核心设计}

WebAssembly 最初面向 Web 应用场景提出，但其设计目标并不局限于浏览器，而是提供一种安全、紧凑、可验证且跨平台的低级二进制表示\cite{haas2017bringing}。从执行模型的核心逻辑来看，Wasm 之所以能兼顾高性能和高安全性，关键在于其建立了栈式虚拟机、线性内存和沙箱化隔离三大核心机制。这些机制既保证了 Wasm 模块能够以统一格式在不同平台上高效运行，也为浏览器内外各类场景中的快速验证、即时编译和便捷部署提供了坚实的底层支撑\cite{mao2024ewapa}。

\subsection{栈式虚拟机}

WebAssembly 采用栈式虚拟机为核心执行机制，其函数代码由一系列作用于隐式操作数栈的指令组成，指令的执行逻辑遵循“从栈顶弹出操作数，执行运算后将结果重新压入栈中”的规则\cite{haas2017bringing}。与寄存器式指令表示相比，栈式表示并不必然带来更高的执行速度，但能通过精简指令结构有效减少字节码体积，使程序表示更加紧凑——这一点对于网络传输和跨平台分发尤为重要。同时，Wasm 采用结构化控制流而非任意跳转，分支目标通过块嵌套的层级关系明确表达，从根本上避免了指令跳入非法位置或形成难以验证的控制流结构。这种设计使得 Wasm 模块能够被单遍验证与单遍编译，降低了运行时的加载成本与实现复杂度，同时为编译优化、性能提升提供了良好的结构基础。

\subsection{线性内存}

Wasm 将程序主存设计为线性内存，即一段连续的字节数组，程序通过字节地址完成对数据的访问与操作\cite{haas2017bringing}。这一设计与传统程序的内存模型较为接近，便于将 C/C++ 等底层编程语言中的指针操作、数组结构及动态内存分配机制高效地映射到 Wasm 执行环境中，降低了原生代码向 Wasm 迁移的成本。在线性内存机制中，程序通过 load/store 指令实现数据读写，内存容量可以按页扩展，从而在保持抽象统一的同时支持较灵活的内存管理。在 Web 场景下，Wasm 线性内存可以与 JavaScript 的 `ArrayBuffer` 实现高效交互；在非 Web 场景下，这种线性地址空间设计也使得 Wasm 运行时更容易与 WASI 等宿主接口对接。更重要的是，Wasm 对所有的内存访问操作都施加严格的边界检查，越界访问将触发 trap，从而有效规避原生代码中常见的越界访问等未定义行为直接破坏执行环境。

\subsection{沙箱设计}

沙箱化执行模型是 Wasm 安全性的核心保障，其核心逻辑在于实现不同执行单元的严格隔离。Wasm 模块的线性内存与运行时内部数据结构、执行栈、本地变量以及其他宿主进程的内存空间相互隔离，模块无法直接访问宿主地址空间中的任意位置。Wasm 模块与外部环境（包括宿主环境、其他模块）的交互必须通过显式导入的宿主函数或标准化系统接口完成，其可访问的资源范围、操作权限均由宿主环境统一管控。这种沙箱化设计使得不可信代码可以和宿主在同一地址空间中安全运行而不破坏宿主环境，为浏览器中的脚本隔离、服务端 Wasm 运行时以及基于 WASI 的系统调用封装提供了基础。也正因如此，WebAssembly 不仅是一种高效的跨平台代码分发格式，更是一种兼顾安全隔离与执行效率的运行时抽象。


\section{WebAssembly 性能研究现状}

目前，学术界与工业界已开展大量相关研究，从多维度对 WebAssembly 的性能特征进行系统分析与研究，研究方向主要涵盖不同运行时性能对比、不同编程语言编译为 Wasm 后的性能差异、Wasm 与 JavaScript 及原生代码（Native Code）的性能对比等核心领域，为 Wasm 性能优化及应用实践提供了重要理论与实验支撑。

\subsection{不同 WebAssembly 运行时性能对比}

随着 WebAssembly 技术的应用场景逐渐从传统浏览器拓展至边缘计算、物联网、服务端等非 Web 领域，不同 Wasm 运行时之间的性能差异逐渐成为研究的热点之一。相关综述性研究指出，当前主流 Wasm 运行时（如 Wasmtime、Wasmer、WAMR、Wasm3、WasmEdge 等）在执行模式（解释执行、即时编译 JIT、提前编译 AOT）、启动开销、吞吐量及内存占用等关键性能指标上存在明显的设计权衡与差异\cite{zhang2025research}。Dierickx 等开展的对比研究进一步证实，不存在一种能够在所有工作负载（workload）上都稳定占优的通用 Wasm 运行时\cite{dierickxcomparative}。

进一步针对 WASI 实现的研究表明，Wasm 性能差异的产生不仅源于 Wasm 字节码本身的执行效率，还与宿主系统接口的实现方式密切相关。eWAPA 研究将 Wasm 程序运行时间进一步拆解为总执行时间、启动时间、WASI 执行时间及系统调用（syscall）时间，研究结果显示，运行时的 WASI 接口实现方案可能成为 I/O 密集型程序性能提升的关键瓶颈\cite{mao2024ewapa}。Jiang 等针对 5 个主流服务端 Wasm 运行时开展差分测试，识别出多个已被开发者确认的性能缺陷，证实运行时内部的实现细节会对服务端场景下 Wasm 的实际性能表现产生显著影响\cite{jiang2023revealing}。上述研究均表明，Wasm 性能表现具有较强的运行时依赖性，运行时的设计与实现是影响 Wasm 性能的关键因素之一。

\subsection{不同编程语言编译为 WebAssembly 性能对比}

另一类研究聚焦于源编程语言、编译工具链及优化选项对 Wasm 性能的影响机制。Prabakar 和 Kiran 对 C++ 与 Rust 两种编程语言编译为 Wasm 后的执行性能进行对比分析，研究结果表明，即便在相同的目标运行平台下，不同编程语言的生态特性及配套工具链仍会引入明显性能差异\cite{prabakar2024webassembly}。

Yan 等开展的系统性研究进一步指出，当前多数 Wasm 编译器都建立在 LLVM 等通用编译基础设施之上，其优化策略并非专门针对 Wasm 的指令集特性与执行模型设计，因此部分在原生代码编译中有效的优化手段，应用于 Wasm 编译时可能无法达到预期效果，甚至出现反直觉的性能退化现象\cite{yan2021understanding}。该研究同时表明，编译器的选择及编译优化级别的设定，不仅会影响 Wasm 代码的执行时间，还会对生成代码的内存占用和跨浏览器兼容性产生影响。综上，现有针对编程语言编译的对比研究揭示，Wasm 性能表现不仅由源编程语言本身决定，还与编译链路和后端优化策略密切相关。


\subsection{基于 Web 的 WebAssembly 与 JavaScript 性能对比}

在浏览器应用场景中，Wasm 与 JavaScript 的性能对比是最早开展、研究成果最丰富的研究方向之一。Sunarto 等对该领域的相关研究进行系统性综述，指出现有研究得出的性能结论并不一致，Wasm 与 JavaScript 的性能表现受浏览器内核、运行设备、应用类型及评价指标等多因素的共同影响\cite{sunarto2023systematic}。Yan 等通过对 41 组程序在多浏览器和多输入规模下的实验发现，Wasm 并不总是稳定优于 JavaScript；JavaScript 的 JIT 优化收益更为显著，而 Wasm 在部分场景中并未因 JIT 优化获得同等幅度提升，且其内存占用通常高于 JavaScript\cite{yan2021understanding}。这一结论说明，Wasm 对 JavaScript 的优势更依赖于负载的计算密度以及初始化开销能否被后续执行摊薄，而不能简单概括为“Wasm 一定更快”。

此外，Wasm 与原生代码（Native Code）的性能关系也并非简单的“接近原生”。Jangda 等基于 SPEC CPU 基准测试套件的研究指出，相比早期基于小型 PolyBench kernel 得出的乐观结论，更大规模应用中 Wasm 相对 native 仍存在显著性能差距，平均执行速度慢约 45\% 到 55\%，其性能损耗主要源于额外的加载/存储（load/store）操作、分支指令（branch）执行及动态安全检查带来的开销\cite{jangda2021not}。Spies 和 Mock 对非浏览器环境下 Wasm 的性能评估也表明，Wasm 在非 Web 场景中是否接近原生性能，仍高度依赖运行时、基准测试类型和系统交互方式\cite{spies2021evaluation}。


总体来看，现有工作已经从浏览器环境内的 JavaScript/Wasm 比较逐渐扩展到 Wasm 与原生代码的性能差异分析、服务端运行时性能评估及 WASI 接口性能优化等方向，为深入理解 WebAssembly 性能特征提供了重要基础。然而，当前多数研究仍以总体运行时间、启动开销或系统调用时间等动态性能指标为研究对象，重点回答“谁更快”以及“运行时性能瓶颈在哪里”，而对程序本身的静态结构如何影响 Wasm 与 native 的性能差异关注不足。尤其在非 Web 场景下，目前仍缺少将程序静态特征与 Wasm 与原生代码的性能差距进行系统关联的研究。基于此，本文从程序静态特征视角出发，对非 Web 场景下 Wasm 与原生代码的性能差距及其结构性来源展开进一步分析与探究。


\section{基于程序静态特征的性能分析相关研究}


基于静态程序特征进行性能建模在系统与编译器领域已有较为深厚的研究基础。与依赖实际执行、硬件计数器或运行时插桩的动态分析方法相比，静态分析能够在不运行程序的前提下，从源码、抽象语法树（AST）或中间表示（IR）中提取程序规模、控制流结构、指令类型和访存模式等信息，具有分析成本低、可迁移性强和可解释性较好的优势。在编译优化研究领域，Dubach 等提出基于代码特征的性能预测方法，利用程序特征估计不同优化方案带来的性能变化，从而降低迭代编译搜索的代价\cite{dubach2007fast}。该研究说明程序自身的结构信息可以为性能建模提供有效的特征信号，为静态特征应用于性能分析提供了理论支撑。

在更直接的静态性能预测研究中，Ardalani 等提出 Static XAPP 框架，仅基于 CPU 源代码的静态分析结果，在不运行程序的情况下预测其跨架构迁移后的加速潜力\cite{ardalani2019static}。该工作在 IR 层面设计了内存访问、分支发散、可并行性以及算术强度等静态特征，并借助机器学习模型对程序性能区间进行分类预测，取得了较高的预测准确率。这一研究结果表明，即使缺少运行时信息，程序在编译中间表示上呈现出的结构属性仍然能够反映其在不同执行平台上的性能行为。这一点对本文研究具有直接的启发意义：Wasm 与 native 虽然运行时环境和代码生成链路不同，但二者之间的性能差异同样可能与计算密度、访存密度、控制流复杂度及函数调用行为等静态结构特征有关。

除人工设计的统计类静态特征外，近年来也有研究尝试直接利用代码结构开展静态性能分析。Pinnow 等将源代码表示为 AST，并采用比较式建模方法预测两个功能正确的程序版本中哪一个运行速度更快\cite{pinnow2021comparative}。该研究指出，绝对执行时间会受到输入规模、系统环境和硬件细节等多种因素影响，单纯依赖静态信息进行绝对性能预测通常较为困难；相较之下，基于代码结构差异的相对性能判断更具可行性。这一思路与本文的研究目标具有一致性。本文并不试图脱离具体实验环境预测程序的绝对运行时间，而是关注同一程序在 native、Wasm JIT 和 Wasm AOT 三种执行模式下的相对性能差异，并进一步分析这些差异与程序静态特征之间的内在关联。

综上，本文采用静态分析方法开展 Wasm 与原生代码的性能解释及预测，在方法论上具有合理性与可行性。一方面，静态特征可在大规模 benchmark 上低成本、统一地提取，便于开展后续统计分析；另一方面，这些特征具有较强可解释性，能够帮助研究者从程序结构角度讨论 Wasm 与 native 性能差距的潜在来源。需要说明的是，本文并不以静态特征替代动态测量，而是将静态分析与实测性能结果结合起来，用于解释性能差异，而非单独进行黑盒式的性能预测。

